\chapter{Use of Artificial Intelligence (AI) Tools}
\label{app:ai}

\section{AI Work Declaration}
\label{app:ai_work}

Three AI tools were used in this project, each for a different purpose:
\begin{itemize}[label=\textbullet]
  \item \textbf{Claude (Anthropic)} was used for logical and complex tasks
    such as troubleshooting, code review, code generation, data-processing
    scripts and planning, and for editing, converting and checking report
    text.
  \item \textbf{GitHub Copilot}, within VS Code, was used for
    autocompletion of spelling and citations while writing.
  \item \textbf{Gemini (Google)} was used for restructuring language in
    the report.
\end{itemize}
I used AI as a tool to support my own work, not to replace it. Every
output was reviewed, tested or checked by me before it was used.

In the project work, AI was used for the following:
\begin{itemize}[label=\textbullet]
  \item \textbf{Troubleshooting.} Diagnosing hardware and firmware problems
    during building and testing, by discussing symptoms and suggesting
    likely causes and checks, which I then carried out on the hardware.
  \item \textbf{Exploring solutions.} Used as an improved search engine,
    suggesting and comparing alternative approaches to design problems.
    I made the choice between them, for the reasons given in the report. 
  \item \textbf{Literature search.} Suggesting relevant papers and search terms for Chapter~2. Every source was located, read and checked by me before it was cited.
  \item \textbf{Code review and improvement.} Restructuring weak or
    repetitive sections of the firmware, checking the code for loopholes
    and edge cases, and adding safety measures such as input limits,
    time-outs and overload handling (Section~\ref{sec:sw_ranging}), as
    well as helping me understand the AD5941 and how its drivers work.
  \item \textbf{Web dashboard code.} The HTML page served by the
    OpenStat's web server was generated by Claude from my description of
    the layout and the data it had to show. I integrated it into the
    firmware and tested it on the device
    (Chapter~\ref{chap:software_design}).
  \item \textbf{Data-comparison scripts.} Help with setting up Python
    scripts to process and compare the data measured on the OpenStat and
    the PalmSens4 (Chapter~\ref{chap:testing}).
  \item \textbf{Planning.} Helping to plan the project schedule
    (Appendix~\ref{app:plan_schedule}) and to break tasks into manageable
    steps, which improved time efficiency.
  \item \textbf{Calibration.} Help setting up the two-point calibration procedure and the calculation of the gain and offset constants for each RTIA range (Section~6.2, Appendix~D.5). The measurements were taken by me, and the resulting constants were checked at the independent $+0.25$~V point before being compiled into the firmware.
\end{itemize}

In the documentation (this report), AI was used for the following:
\begin{itemize}[label=\textbullet]
  \item \textbf{Firmware chapter.} For Chapter~\ref{chap:software_design}
    (Firmware Design), Claude was given my rough drafts of the flow
    diagrams and translated them into \LaTeX{} (TikZ). I checked the
    diagrams against the code and the behaviour of the device, and
    corrected them where they did not match.
  \item \textbf{Derivations.} Claude was used to simplify and clarify the
    derivations in Chapter~\ref{chap:modeling} and \ref{chap:hardware_design}. I checked each step of
    the resulting derivations myself.
  \item \textbf{Drafting.} Text in Chapter~\ref{chap:testing} (Testing and
    Calibration), Chapter~\ref{chap:environmental} (Environmental and
    Sustainability Analysis), Chapter~\ref{chap:conclusion} (Conclusion
    and Future Work) and Appendix~\ref{app:outcomes} (Graduate Attributes)
    was improved with Claude, which flagged incomplete sentences and
    inconsistencies. I reviewed, corrected, edited and checked the content.

  \item \textbf{Graduate attributes.} Claude was used to help identify
    what each graduate attribute required. The evidence and discussion in
    Appendix~\ref{app:outcomes} are drawn from my own work.
  \item \textbf{Conversion to \LaTeX{}.} Chapter~\ref{chap:hardware_design}
    (Hardware Design) was written by me as a rough draft in Word and
    converted into \LaTeX{} with Claude, which also flagged incomplete
    sentences and inconsistencies for me to resolve.
  \item \textbf{Bill of materials.} Claude compiled the bill of materials
    for this report from the components I selected.
  \item \textbf{Language editing.} AI was used in Chapters~\ref{chap:introduction}
    to~\ref{chap:modeling} only to improve structure, clarity and flow,
    mainly with Gemini.
  \item \textbf{\LaTeX{} support.} Finding and fixing compilation errors,
    and checking that all figures, tables, equations and sections are
    correctly referenced.
  \item \textbf{Proofreading.} Checking spelling, with Copilot also
    completing citations while writing in VS Code.
\end{itemize}

All AI outputs were verified before use. Firmware changes were compiled,
flashed and run on the OpenStat, and their behaviour was checked against
the serial output. An oscilloscope was used to confirm that the code
produced the correct DPV excitation waveform (Section~\ref{sec:waveform}).
Not all AI-suggested code worked, some aspects failed and had to be debugged and corrected by me before they
behaved as intended.

AI was not used to generate or alter any measured data. All results and
conclusions in this report are based on measurements made during testing.

\section{Own Work Declaration}
\label{app:own_work}

In the project work, the following was done by me:
\begin{itemize}[label=\textbullet]
  \item \textbf{Project definition.} Defining the aim, scope and
    requirements of the project with the supervisor
    (Chapter~\ref{chap:introduction}, Tables~4.1 and~4.2).
  \item \textbf{Literature review.} Selecting the
    sources reviewed in Chapter~2.
  \item \textbf{Hardware design.} Selecting the components, designing the
    schematic and PCB layout in KiCad, ordering the board, and assembling
    and bringing up the prototype (Chapter~\ref{chap:hardware_design}).
  \item \textbf{Firmware.} Writing the firmware, designing the
    code structure, including current measurement and range selection,
    and making the design decisions behind them
    (Chapter~\ref{chap:software_design}). Except the content in \ref{app:ai_work}
  \item \textbf{User interface.} Designing the layout of the onboard
    screens and web dashboard, and how the user interacts with the device.
  \item \textbf{Testing and calibration.} Carrying out all measurements,
    including the waveform, resistor calibration, Randles circuit, salt
    and ferricyanide tests, preparing the test solutions and operating
    the PalmSens4 (Chapter~\ref{chap:testing}).
  \item \textbf{Analysis and conclusions.} Interpreting the results and
    drawing the conclusions of Chapters~\ref{chap:testing}
    to~\ref{chap:conclusion}.
\end{itemize}

In the documentation (this report), the following was done by me:
\begin{itemize}[label=\textbullet]
  \item \textbf{Writing.} Writing Chapters~\ref{chap:introduction}
    to~\ref{chap:modeling}, including the original derivations, the rough
    draft of Chapter~\ref{chap:hardware_design}, the rough flow diagrams
    of Chapter~\ref{chap:software_design}, and the rough drafts and notes
    used for Chapters~\ref{chap:testing} to~\ref{chap:conclusion} and
    Appendix~\ref{app:outcomes}.
  \item \textbf{Content.} Deciding what the report covers, and selecting
    the results, figures and tables presented in it.
  \item \textbf{Review.} Reviewing, correcting and approving the final
    text, figures, derivations and references of the whole report.
\end{itemize}